\section{Conclusion and Discussion}
\label{sec:conclusion_discussion}

This paper identifies \emph{memory laundering} as a safety blind spot in memory-augmented LLM agent systems. The failure mode is not that compressed memories remain overtly toxic. Rather, summarization can remove the lexical markers detected by toxicity classifiers while preserving enough adversarial framing to steer downstream agents. A memory state can therefore be classifier-clean under standard thresholds yet remain behaviorally influential when reused as context.

To measure this hidden regime, we introduced the sub-threshold propagation gap (SPG), which quantifies downstream toxicity differences conditioned only on memory states that a deployed monitor would label safe. SPG complements the paired propagation gap $\Delta\mu$: $\Delta\mu$ captures overt propagation through the transcript, while SPG captures hidden propagation through compressed memory. Together, these metrics show that transcript and memory are distinct state channels rather than redundant views of the same phenomenon.

The mitigation results support a systems-level view of agent safety. Output filtering treats toxicity as a local property of a generated message and therefore misses contamination that survives in state. DPO reduces parametric amplification but does not prevent toxic framing from re-entering future contexts through transcript or memory. Memory-only sanitization helps on the hidden channel but does not close transcript backflow. The full state-aware system performs best because it combines read-side context sanitization, write-side state control, and parameter-level debiasing.

The broader implication is that safety for stateful agents cannot be solved by model alignment alone. Agents repeatedly observe, generate, summarize, and update shared state; failures can therefore arise from the interaction between model behavior and state dynamics. Robust evaluation should include counterfactual multi-turn rollouts, explicit channel decomposition, and metrics that test whether apparently safe intermediate states still influence downstream behavior. Robust mitigation should intervene before contaminated content is compressed into persistent memory. In short: for memory-augmented agents, safety requires sanitizing the state before summarizing it.
